% ============================================================
%  MTH636 – Game Theory | Lecture 2
%  Outcomes, Preferences, Utility, and Normal Form Games
% ============================================================
\documentclass[12pt, a4paper]{article}

% ============================================================
%  MTH636 -- Game Theory
%  Notes Preamble — Color-Coded, TikZ-Enhanced
% ============================================================

% --- Core packages ---
\usepackage[T1]{fontenc}
\usepackage[utf8]{inputenc}
\usepackage{lmodern}
\usepackage[margin=2.5cm, top=2.8cm, bottom=2.8cm]{geometry}
\usepackage{microtype}

% --- Math ---
\usepackage{amsmath, amssymb, amsthm, mathtools}
\usepackage{bm}        % bold math
\usepackage{cancel}    % strike-through in math

% --- Color & Boxes ---
\usepackage[dvipsnames, svgnames, x11names, table]{xcolor}
\usepackage[most]{tcolorbox}
\tcbuselibrary{breakable, skins, theorems}

% --- TikZ & Diagrams ---
\usepackage{tikz}
\usepackage{tikz-qtree}
\usetikzlibrary{arrows.meta, positioning, shapes.geometric, fit, backgrounds,
                decorations.pathreplacing, calc, matrix, trees}
\usepackage{pgfplots}
\pgfplotsset{compat=1.18}

% --- Tables ---
\usepackage{booktabs}
\usepackage{multirow}
\usepackage{array}
\usepackage{tabularx}
\usepackage{makecell}

% --- Typography & Layout ---
\usepackage{enumitem}
\usepackage{parskip}
\usepackage{fancyhdr}
\usepackage{titlesec}
\usepackage{titletoc}
\usepackage{hyperref}
\usepackage{graphicx}
\usepackage{float}
\usepackage{caption}

% --- Color Palette ---
\definecolor{defblue}{HTML}{1A5276}      % Definition header
\definecolor{defbluebg}{HTML}{D6EAF8}    % Definition background
\definecolor{thmgreen}{HTML}{1E8449}     % Theorem header
\definecolor{thmbg}{HTML}{D5F5E3}        % Theorem background
\definecolor{exorange}{HTML}{B7770D}     % Example header
\definecolor{exbg}{HTML}{FEF9E7}         % Example background
\definecolor{warnred}{HTML}{922B21}      % Warning/Key Insight header
\definecolor{warnbg}{HTML}{FADBD8}       % Warning background
\definecolor{noticegray}{HTML}{2C3E50}   % Note header
\definecolor{noticebg}{HTML}{EBF5FB}     % Note background
\definecolor{taketawayyellow}{HTML}{7D6608} % Key Takeaway header
\definecolor{takeawaybg}{HTML}{FDFEFE}   % Takeaway background
\definecolor{accentpurple}{HTML}{6C3483} % Accent / proof
\definecolor{lightgray}{HTML}{F2F3F4}    % shading

% ============================================================
%  TCOLORBOX ENVIRONMENTS
% ============================================================

% -- Definition (Blue) --
\newtcolorbox{defbox}[2][]{
  enhanced, breakable,
  colback=defbluebg, colframe=defblue, coltitle=white,
  fonttitle=\bfseries\sffamily,
  title={Definition: #2},
  attach boxed title to top left={yshift=-2mm, xshift=4mm},
  boxed title style={colback=defblue, rounded corners},
  arc=4pt, boxrule=1.2pt,
  left=6pt, right=6pt, top=8pt, bottom=6pt,
  #1
}

% -- Theorem / Proposition (Green) --
\newtcolorbox{thmbox}[2][]{
  enhanced, breakable,
  colback=thmbg, colframe=thmgreen, coltitle=white,
  fonttitle=\bfseries\sffamily,
  title={Theorem: #2},
  attach boxed title to top left={yshift=-2mm, xshift=4mm},
  boxed title style={colback=thmgreen, rounded corners},
  arc=4pt, boxrule=1.2pt,
  left=6pt, right=6pt, top=8pt, bottom=6pt,
  #1
}

% -- Proposition (Green variant) --
\newtcolorbox{propbox}[2][]{
  enhanced, breakable,
  colback=thmbg, colframe=thmgreen!70!black, coltitle=white,
  fonttitle=\bfseries\sffamily,
  title={Proposition: #2},
  attach boxed title to top left={yshift=-2mm, xshift=4mm},
  boxed title style={colback=thmgreen!70!black, rounded corners},
  arc=4pt, boxrule=1.2pt,
  left=6pt, right=6pt, top=8pt, bottom=6pt,
  #1
}

% -- Example (Orange/Yellow) --
\newtcolorbox{exbox}[2][]{
  enhanced, breakable,
  colback=exbg, colframe=exorange, coltitle=white,
  fonttitle=\bfseries\sffamily,
  title={Example: #2},
  attach boxed title to top left={yshift=-2mm, xshift=4mm},
  boxed title style={colback=exorange, rounded corners},
  arc=4pt, boxrule=1.2pt,
  left=6pt, right=6pt, top=8pt, bottom=6pt,
  #1
}

% -- Key Insight / Warning (Red) --
\newtcolorbox{warnbox}[2][]{
  enhanced, breakable,
  colback=warnbg, colframe=warnred, coltitle=white,
  fonttitle=\bfseries\sffamily,
  title={Key Insight: #2},
  attach boxed title to top left={yshift=-2mm, xshift=4mm},
  boxed title style={colback=warnred, rounded corners},
  arc=4pt, boxrule=1.2pt,
  left=6pt, right=6pt, top=8pt, bottom=6pt,
  #1
}

% -- Note (Gray-Blue) --
\newtcolorbox{notebox}[1][]{
  enhanced, breakable,
  colback=noticebg, colframe=noticegray!60, coltitle=white,
  fonttitle=\bfseries\sffamily,
  title={Note},
  attach boxed title to top left={yshift=-2mm, xshift=4mm},
  boxed title style={colback=noticegray!60, rounded corners},
  arc=4pt, boxrule=0.8pt,
  left=6pt, right=6pt, top=8pt, bottom=6pt,
  #1
}

% -- Key Takeaway --
\newtcolorbox{takebox}[1][]{
  enhanced, breakable,
  colback=takeawaybg, colframe=taketawayyellow, coltitle=white,
  fonttitle=\bfseries\sffamily,
  title={Key Takeaway},
  attach boxed title to top left={yshift=-2mm, xshift=4mm},
  boxed title style={colback=taketawayyellow, rounded corners},
  arc=4pt, boxrule=1.2pt,
  left=6pt, right=6pt, top=8pt, bottom=6pt,
  #1
}

% -- Intuition --
\newtcolorbox{intubox}[1][]{
  enhanced, breakable,
  colback=accentpurple!8, colframe=accentpurple!60, coltitle=white,
  fonttitle=\bfseries\sffamily,
  title={Intuition},
  attach boxed title to top left={yshift=-2mm, xshift=4mm},
  boxed title style={colback=accentpurple!80, rounded corners},
  arc=4pt, boxrule=0.8pt,
  left=6pt, right=6pt, top=8pt, bottom=6pt,
  #1
}

% ============================================================
%  PAYOFF MATRIX HELPER
%  Usage: inside tabular, use \payoff{a}{b} for (a,b)
% ============================================================
\newcommand{\payoff}[2]{$(#1,\;#2)$}
\newcommand{\nash}[2]{\underline{$#1$},\underline{$#2$}}  % underline NE payoffs
\newcommand{\Nset}{\mathbb{N}}
\newcommand{\Rset}{\mathbb{R}}
\newcommand{\Sset}{\mathcal{S}}
\newcommand{\Ui}{u_i}

% ============================================================
%  SECTION FORMATTING
% ============================================================
\titleformat{\section}
  {\large\bfseries\sffamily\color{defblue}}
  {\thesection}{1em}{}[\titlerule]

\titleformat{\subsection}
  {\normalsize\bfseries\sffamily\color{thmgreen!80!black}}
  {\thesubsection}{1em}{}

\titleformat{\subsubsection}
  {\normalsize\itshape\sffamily\color{exorange!80!black}}
  {\thesubsubsection}{1em}{}

% ============================================================
%  HEADER / FOOTER
% ============================================================
\pagestyle{fancy}
\fancyhf{}
\fancyhead[L]{\sffamily\small\color{defblue} MTH636 -- Game Theory}
\fancyhead[R]{\sffamily\small\color{defblue} \leftmark}
\fancyfoot[C]{\sffamily\small\thepage}
\renewcommand{\headrulewidth}{0.4pt}

% ============================================================
%  HYPERREF SETUP
% ============================================================
\hypersetup{
  colorlinks=true,
  linkcolor=defblue,
  urlcolor=thmgreen,
  citecolor=accentpurple,
  bookmarksnumbered=true,
  pdftitle={MTH636 Game Theory Notes},
  pdfauthor={MTH636 Student Notes}
}

% ============================================================
%  MISC HELPERS
% ============================================================
\newcommand{\BR}{\text{BR}}        % Best Response
\newcommand{\NE}{\text{NE}}        % Nash Equilibrium
\newcommand{\SPNE}{\text{SPNE}}    % Subgame Perfect NE
\newcommand{\Dom}{\text{dom}}      % Dominance
\newcommand{\mixs}[1]{\sigma_{#1}} % Mixed strategy
\newcommand{\pures}[1]{s_{#1}}     % Pure strategy


% ============================================================
\begin{document}

% Title
\begin{center}
  {\LARGE\bfseries\sffamily\color{defblue} MTH636 -- Game Theory}\\[6pt]
  {\Large\sffamily Lecture 2: Preferences, Utility, and Normal Form Games}\\[4pt]
\end{center}
\vspace{0.5cm}
\hrule height 1.5pt
\vspace{0.3cm}

% ============================================================
\section{Revisiting the Beauty Contest Game}

Recall the Beauty Contest Game with $N = \{1, 2, 3\}$. Each player chooses a number between $0$ and $100$. The target is $\frac{2}{3}$ of the average, and the player whose choice is closest to the target wins.

The set of possible \textbf{outcomes} (denoted by $O$) includes:
\begin{itemize}
    \item Player 1 is the sole winner.
    \item Player 2 is the sole winner.
    \item Player 3 is the sole winner.
    \item Players 1 and 2 are joint winners.
    \item Players 1 and 3 are joint winners.
    \item Players 2 and 3 are joint winners.
    \item All are winners.
\end{itemize}

\subsection{Preferences over Outcomes}
Player 1 has a preference relation over these outcomes. Naturally, Player 1 prefers to win alone rather than share the win, and prefers to share the win rather than not win at all. We denote this preference relation by $\succsim_1$.

\begin{itemize}
    \item Player 1 is the sole winner $\succ_1$ Players 1 and 2 are joint winners $\sim_1$ Players 1 and 3 are joint winners $\succ_1$ All are winners $\succ_1$ all the other outcomes.
\end{itemize}

Formally, a preference relation $\succsim_i$ is a subset of $O \times O$. When we write $(x, y) \in {\succsim_i}$, or simply $x \succsim_i y$, it means Player $i$ \textbf{weakly prefers} outcome $x$ over outcome $y$.

% ============================================================
\section{Utility Functions}

\begin{defbox}{Utility Function}
  Suppose $O$ is the set of outcomes. Then a \textbf{utility function} $u$ is a mapping from $O$ to $\mathbb{R}$.
\end{defbox}

\begin{defbox}{Representation of a Preference}
  A utility function $u$ \textbf{represents} a preference relation $\succsim$ over $O$ if for all $x, y \in O$:
  \[
    x \succsim y \iff u(x) \geq u(y)
  \]
\end{defbox}

\begin{thmbox}{Existence of a Utility Representation}
  Suppose $O$ is finite. Then a preference relation $\succsim$ can be represented by a utility function if and only if the preference is \textbf{complete} and \textbf{transitive}.
  
  \textit{Note: If $O$ is infinite, a similar result holds but requires additional continuity assumptions.}
\end{thmbox}

% ============================================================
\section{Normal (Strategic) Form Games}

\begin{defbox}{Normal Form Game}
  A normal form (or strategic form) game is a 3-tuple $\langle N, (S_i)_{i \in N}, (u_i)_{i \in N} \rangle$ where:
  \begin{itemize}
      \item $N$ is a finite set of players.
      \item $S_i$ is the set of strategies available to player $i$.
      \item $u_i : \prod_{j \in N} S_j \to \mathbb{R}$ is the utility function for player $i$, which maps a strategy profile to a real-valued payoff.
  \end{itemize}
\end{defbox}

\subsection{Formalizing the Beauty Contest Game}
Using this definition, we can formally define the Beauty Contest Game:
\begin{itemize}
    \item $N = \{1, 2, 3\}$
    \item $S_1 = [0, 100]$, $S_2 = [0, 100]$, $S_3 = [0, 100]$
    \item Utility functions:
    \[
      u_1(x_1, x_2, x_3) = \begin{cases} 
        \frac{1}{|W(x_1, x_2, x_3)|} & \text{if } 1 \in W(x_1, x_2, x_3) \\ 
        0 & \text{otherwise} 
      \end{cases}
    \]
    where $W(x_1, x_2, x_3) = \arg\min_{i \in N} \left|x_i - \frac{2}{3}\bar{x}\right|$ is the set of winning players.
\end{itemize}

% ============================================================
\section{Classic 2x2 Games}

\begin{exbox}{Prisoner's Dilemma}
  $N = \{P_1, P_2\}$, $S_1 = S_2 = \{C, D\}$ (Cooperate, Defect)
  
  \begin{center}
  \renewcommand{\arraystretch}{1.4}
  \begin{tabular}{c|cc}
    $P_1 \setminus P_2$ & C & D \\
    \hline
    C & \payoff{-1}{-1} & \payoff{-5}{0} \\
    D & \payoff{0}{-5} & \payoff{-3}{-3} \\
  \end{tabular}
  \end{center}
  
  In this game, players will play $(D, D)$ because for both of them, choosing $D$ strictly dominates choosing $C$.
\end{exbox}

\begin{exbox}{Battle of the Sexes}
  $N = \{\text{Man}, \text{Woman}\}$, $S_M = S_W = \{F, D\}$ (Football, Dance)
  
  \begin{center}
  \renewcommand{\arraystretch}{1.4}
  \begin{tabular}{c|cc}
    $M \setminus W$ & F & D \\
    \hline
    F & \payoff{2}{1} & \payoff{0}{0} \\
    D & \payoff{0}{0} & \payoff{1}{2} \\
  \end{tabular}
  \end{center}
  
  This is a classic coordination game where players wish to coordinate on the same event but have differing preferences over which event to attend.
\end{exbox}

\end{document}
